\section{Related Work}
\label{sec:related}

\paragraph{Temporal dynamics of the web.}
The rate at which web pages change has been studied since the early days of web
crawling. \citet{cho2003estimating} showed that page change frequency follows a
Poisson process, and that a crawler with a fixed budget should prioritize
high-churn pages. \citet{fetterly2003evolution} measured change rates for a
large sample of web pages and found that roughly 40\% of pages change within a
week, with significant variance across domains. \citet{ntoulas2004s} studied the
evolution of web content over time and found that new pages are created at a rate
far exceeding the capacity of any single crawler. FreshState builds on this
tradition but focuses on \emph{answer-bearing spans}---the specific values in a
page that determine whether a search result is factually accurate---rather than
coarse page-level change signals.

\paragraph{Web archiving and temporal snapshots.}
The Internet Archive's Wayback Machine~\citep{internetarchive} maintains a crawl
history for billions of URLs, enabling retrospective study of how pages changed
over time. The CDX API provides programmatic access to this history at URL
granularity~\citep{cdxapi}. Temporal IR systems have exploited Wayback snapshots
to study document freshness~\citep{joho2013temporalir}, and systems like
TWIX~\citep{whiting2013twix} query the archive to reconstruct past page states.
Our retrospective strategy (Strategy A) leverages the CDX API to collect
historically stable candidate URLs, while Strategy B bypasses the archive entirely
in favor of live monitoring. This distinction matters: Wayback crawls are sparse and
irregular for many domains (the CDX API was unreachable for several hours during our
data collection, illustrating its reliability limitations), whereas prospective
monitoring provides complete, evenly-spaced snapshots.

\paragraph{Temporal fact verification.}
The claim verification literature has begun to address temporal scope.
TempLAMA~\citep{dhingra2022templama} probes whether language models retain
time-sensitive relational facts, and CREAK~\citep{onoe2021creak} includes commonsense
claims that are true at one point in time and false at another; both rely on manually
annotated claims rather than automatically extracted change events. \citet{chen2021timeqa}
introduced time-sensitive QA but used Wikipedia revision history as the ground
truth, which captures only deliberate human edits---not the kind of silent,
real-world changes (a listed apartment renting out, a product going out of stock)
that FreshState targets. \citet{kasner2023tabgenie} and \citet{shi2024freshllm}
study knowledge freshness in language models, but treat model knowledge as the
source of staleness rather than search results. FreshState is complementary:
it provides ground-truth change labels at the web-document level, enabling
downstream evaluation of both retrieval systems and language model knowledge.

\paragraph{Retrieval-augmented generation and freshness.}
Retrieval-augmented generation (RAG)~\citep{lewis2020rag} grounds language model
outputs in retrieved documents, making document freshness a first-order concern.
\citet{shi2024freshllm} construct FreshQA, a benchmark of factual questions whose
answers change over time, and evaluate both closed-book LLMs and search-augmented
systems. Their findings show that all evaluated systems degrade significantly on
questions with recently changed answers. FreshState differs from FreshQA in two
key ways: (1) it targets document-level staleness rather than model-level knowledge,
and (2) it uses automatically extracted change events rather than human-authored
question-answer pairs. This makes FreshState cheaper to update and applicable to
domains (rental prices, product availability) that are too ephemeral to support
reliable human annotation.

\paragraph{Search result quality evaluation.}
Evaluation of web search quality has a long history in IR~\citep{voorhees2005trec}.
Most evaluation frameworks measure topical relevance~\citep{craswell2021overview}
or answer correctness~\citep{kwiatkowski2019nq} at the time of annotation. Temporal
freshness has been studied as a secondary quality dimension~\citep{adar2009web}, but
it is rarely the primary evaluation target. The FEVER shared task~\citep{thorne2018fever}
tests fact verification against Wikipedia but does not model temporal change.
FreshState provides the first benchmark specifically designed to evaluate whether
a retrieval system returns a search result whose answer-bearing span is \emph{still
accurate} at query time, as opposed to accurate only at annotation time.
